% SPDX-FileCopyrightText: 2026 Will Cook
% SPDX-License-Identifier: CC-BY-4.0
\documentclass[10pt]{article}
\usepackage{paper-house-style}
\usepackage{xurl}
\usepackage[colorlinks=true,linkcolor=linkinternal,urlcolor=linkexternal,citecolor=linkinternal]{hyperref}
\hypersetup{pdftitle={Writing a Good Mathematical Paper},pdfauthor={Will Cook},pdflang={en-GB},bookmarksnumbered=true}
\runninghead{Writing a good mathematical paper}
\title{Writing a Good Mathematical Paper}
\subtitle{A compact guide for authors and AI assistants}
\author{Will Cook}
\date{7 October 2026}
\emergencystretch=1.5em
\newcommand{\instruction}[2]{\par\medskip\noindent\textbf{#1.} #2\par}
\begin{document}
\maketitle
\paperprovenance
\papersectiontarget{compact-writing-guide}
Test the draft for its intended reader: state the contribution, reconstruct
the decisive inference and apply the result. Begin where one of these tasks
fails; use the relevant instruction, not every instruction.

\instruction{1. Decide what the paper contributes}{
Choose a coherent principal result; say what it enables or clarifies.
Compared with the nearest prior work, what changes in the assumptions or
conclusion? State what remains open. For a conditional result, give familiar
cases its hypothesis admits or excludes, where known, and show where the proof
uses it. Distinguish contribution from proof mechanism; the abstract needs
both. Neither theorem count nor generality establishes importance.}

\instruction{2. Read the kind of mathematics you are writing}{
Read nearby papers for terminology, hypotheses and attribution; read complete
local arguments for explanation. Match the task: a parameter choice, limit or
estimate. Check the version and passage. Identify what justifies the source's
next step and the corresponding fact in your proof; explain it in your own words.}

\instruction{3. Check the statement before polishing its prose}{
Check domains, quantifiers, signs, endpoints, constants and both directions
of claimed equivalences. State what is fixed and what later choices may depend
on. Keep assumptions with the claim; compare titles, abstracts and conclusions
with it. When properties are needed together, prove they hold at the same
witness: the even and odd integers are both unbounded but disjoint.
Refer a discovered proof gap for mathematical review.}

\instruction{4. Organise by logical dependence}{
Give the question, result and proof idea before technical detail accumulates.
For each major step, say what it uses, what it proves and where that conclusion
is needed. Introduce prerequisites when their purpose is visible; use examples
to explain unfamiliar operations. Let headings name objects, results or
operations; make assumptions and proofs findable from the theorem. Logical
order need not reproduce discovery order.}

\instruction{5. Explain the hard transition}{
Work backwards from the next conclusion: what must be controlled, and how
accurately? State the requirement before estimating or choosing a parameter.
To certify $x\in(1/2,1)$, place an enclosure wholly inside that interval;
$[1/2,1]$ does not suffice. For a small-integer contradiction, justify
integrality, nonvanishing and $|N|<1$ separately. Explain which fact supplies
each requirement; a repeated technique may have a different job.}

\instruction{6. Separate construction from completion}{
Justify each choice and preserved condition, then prove attainment.
For example, $r_n=x-S_n$, $0\le r_n\le\varepsilon_n$ and
$\varepsilon_n\to0$ give $S_n\to x$; repeated choice alone does not.
A classification needs construction and exhaustiveness; an enclosure alone
need not establish attainment. State what stays fixed; justify convergence of
the whole summand, including growing products. To interchange limit and sum,
check the theorem's hypotheses; dominated convergence needs a summable majorant
independent of the limiting parameter.}

% Keep the final reading paragraph with the fourteen instructions.
\newpage\enlargethispage{3\baselineskip}
\instruction{7. Make examples and figures do work}{
Choose a small instance that exposes a definition, invariant, obstruction or
necessary hypothesis; a simpler case may hide it. Show what changes and what
is preserved. When reindexing a sum, check multiplicities, not just values.
Include relevant endpoints and say what remains unproved. A finite check
supports an infinite claim only through a proved reduction. Check that a
witness belongs to the claimed object. Label a diagram's objects, maps and
conventions; keep its explanation nearby.}

\instruction{8. Use the field's words and purposeful notation}{
Keep established terms when their definitions fit, such as ``Stieltjes moment
sequence''. Replace private labels by the actual condition; $u_n/2^n\to0$
is weaker than subexponential growth. Define symbols by their role before use
and keep that role stable. Remove abbreviations whose translation costs more
than their reuse saves. Translate borrowed indices and normalisations.
Distinguish a sum, its partial sums and its remainders, and a value from its enclosure.}

\instruction{9. Make sentences express the argument}{
Replace ``We proceed with the estimate'' by its reason: ``The terms are
nonnegative, so dropping the ordering restriction gives an upper bound.'' Let ``since'' supply a reason
and ``hence'' mark a deduction. Give each ``this'' an unmistakable referent.
After ``similarly'', check that the same hypotheses apply. Punctuate displays
as parts of sentences. Keep parallel claims parallel and conditions beside
their consequences. Remove filler, not mathematical qualifications.}

\instruction{10. Credit the exact input at its use}{
Distinguish applied theorems, adapted constructions, historical attribution,
terminology and expository models. Check hypotheses, version
and passage; identify the input used here. State what an adaptation retains and
changes. Credit supplied ideas and corrections locally. A writing analogy
supplies no proof. Report only reading and checks performed.}

\instruction{11. State what has actually been established}{
Separate scope from evidence: conditional implications may have ordinary or
formal proofs. Match formal support to the exact statement, definitions and
hypotheses. Check coverage of the steps joining lemmas. Missing formalisation
does not invalidate an ordinary proof. Keep unproved inputs beside dependent
results. Label computations and conjectures; give computations' finite scope
locally and reproduction details once. For one-sided tests, failure may be inconclusive.
Detecting each fixed input need not give a uniform stopping bound or the
required witnesses.}

\instruction{12. Preserve useful detail and informative failures}{
Keep omitted detail recoverable by section or theorem references; check their
destinations and reconcile statements, assumptions and evidence in both directions.
For a failed route, retain the attempted implication, assumptions, witness and
missing ingredient. Name the assertion a counterexample refutes; a failed method
need not refute the theorem. Put build commands and source identifiers in the
verification record. Do not invent discovery history.}

\instruction{13. Review changes as claims, then as prose}{
Compare meanings before fluency, including words beside unchanged formulas.
Check earlier criticisms against the current source. Correct overclaims and
stale underclaims without losing hypotheses or credit. Accept, repair, reject
or defer each proposal for a stated reason; preserve useful rejected revisions.
Test a local improvement's limits before promoting it to general advice.}

\instruction{14. Read the delivered paper from a cold start}{
Render the final sources; inspect formulas, captions, page breaks, bibliography
and reference destinations. Read without supplying steps from memory. Can the
reader reconstruct the hard inference, apply the result and identify an
excluded case, where known? Where reconstruction stops, record what is
established and what the next step needs. A successful build is not a comprehension test.
Revise that transition; preserve what works.}

\par\medskip\noindent\small
\textit{Further reading.}
Halmos, \href{https://doi.org/10.5169/seals-43857}{\emph{How to write mathematics}}, \S\S3--5 (1970);
Knuth et al., \href{https://cs.stanford.edu/~knuth/klr.html}{\emph{Mathematical Writing}}, \S\S1--3 (1988);
Gowers, \href{https://gowers.wordpress.com/2007/10/19/my-favourite-pedagogical-principle-examples-first/}{\emph{Examples first!}} (2007);
Tao, \href{https://terrytao.wordpress.com/advice-on-writing-papers/}{\emph{On writing}}.
The \href{https://github.com/wcook04/plectis-erdos/blob/main/paper/exposition/writing-mathematics-from-reviewed-revisions.pdf}{\emph{long companion}}
gives worked cases (\S\S2--4), revision practice (\S5) and genre transfer (\S6).
\end{document}
